
Netdata se caracteriza por tener un consumo de recursos relativamente bajo. Sin embargo, sus requerimientos pueden variar según la cantidad de dispositivos, servicios y métricas que se deseen supervisar, así como el periodo durante el cual se almacenen los datos históricos.

\section{Requerimientos de hardware}

Para ejecutar un agente de Netdata de manera independiente se deben considerar los siguientes recursos:

\begin{itemize}
    \item \textbf{Procesador:} al menos un núcleo de procesamiento. Con la configuración predeterminada, Netdata normalmente utiliza entre el 1\,\% y el 5\,\% de un núcleo. Este consumo puede aumentar al habilitar más recolectores, análisis de anomalías o una mayor cantidad de equipos supervisados.
    
    \item \textbf{Memoria RAM:} aproximadamente entre 100 y 200~MB para una instalación básica. En un entorno de producción típico, el consumo puede encontrarse entre 250 y 350~MB, dependiendo de la cantidad de métricas recopiladas.
    
    \item \textbf{Almacenamiento:} aproximadamente 4~GB con la configuración predeterminada. De este espacio, alrededor de 3~GB se utilizan para almacenar métricas y el espacio restante corresponde a metadatos, alertas y otros archivos internos.
    
    \item \textbf{Interfaz de red:} conexión Ethernet o inalámbrica para permitir el acceso al panel de monitoreo y la comunicación con los dispositivos supervisados. \cite{netdata-recursos} \cite{netdata-almacenamiento}
\end{itemize}

\section{Requerimientos de software}

Para instalar y utilizar Netdata en la plataforma seleccionada se requiere lo siguiente:

\begin{itemize}
    \item Un sistema operativo basado en Linux y compatible con la arquitectura del procesador.
    
    \item Una distribución de 64 bits para arquitectura ARM64 o AArch64, correspondiente al procesador utilizado por la Raspberry Pi 5.
    
    \item Acceso a una cuenta con privilegios administrativos mediante \texttt{sudo}, debido a que la instalación requiere modificar archivos del sistema y crear el servicio de Netdata.
    
    \item Una conexión a Internet durante la instalación, necesaria para descargar Netdata, sus paquetes y las dependencias correspondientes.
    
    \item Las herramientas \texttt{curl} o \texttt{wget}, utilizadas para descargar el instalador oficial.
    
    \item El administrador de servicios \texttt{systemd}, empleado para iniciar Netdata automáticamente junto con el sistema operativo y controlar su estado.
    
    \item Un navegador web moderno, como Mozilla Firefox, Google Chrome o Microsoft Edge, para acceder al panel de monitoreo.
    
    \item Disponibilidad del puerto TCP \texttt{19999}, utilizado de manera predeterminada por la interfaz web local de Netdata. \cite{netdata-seguridad}
\end{itemize}

Netdata incluye su propio motor de almacenamiento de métricas, denominado \texttt{dbengine}. Por esta razón, una instalación básica no requiere desplegar previamente un sistema gestor de bases de datos externo, como MySQL, MariaDB o PostgreSQL. El instalador oficial también se encarga de detectar el sistema operativo y seleccionar el método de instalación más apropiado, utilizando paquetes nativos cuando estos se encuentran disponibles. \cite{netdata-arquitecturas}

\section{Recursos disponibles para el proyecto}

Para el desarrollo del proyecto se dispone del siguiente equipo:

\begin{itemize}
    \item Raspberry Pi 5.
    \item Procesador ARM de cuatro núcleos y 64 bits.
    \item 8~GB de memoria RAM.
    \item Unidad SSD externa con capacidad de 120~GB.
    \item Conexión de red mediante Ethernet.
    \item Ubuntu Server de 64 bits como sistema operativo.
\end{itemize}

El equipo seleccionado supera ampliamente los requerimientos de una instalación individual de Netdata. Los 8~GB de memoria RAM permiten ejecutar el agente, el sistema operativo y servicios adicionales sin comprometer el funcionamiento general. Asimismo, el SSD de 120~GB proporciona suficiente espacio para el programa, sus registros y el almacenamiento histórico de métricas.

El uso de una unidad SSD también ofrece una mayor velocidad de lectura y escritura y una mejor resistencia ante operaciones frecuentes que una tarjeta microSD. Esto resulta conveniente debido a que Netdata almacena continuamente información de monitoreo en su base de datos interna.